\chapter{The erasure specification: erasability, Erases, ErasesEnv}
\label{chap:erases}

This chapter says what it \emph{means} for a \lbox{} term to be a correct erasure of a Lean
term, and proves the metatheory that the forward simulation of
Chapter~\ref{chap:correctness} and the eraser bridge of Chapter~\ref{chap:refinement} spend.
It follows [S, Sect. 7.3]: first a predicate saying which source terms are computationally
irrelevant (\code{Erasable}, [S]'s \code{isErasable}, p. 8:61, and Letouzey's ``logical part
or type scheme'', [L, Def. 1 and Def. 3]); then a \emph{relation} rather than a function,
because erasure is not closed under evaluation ([S, pp. 8:61--8:62]); then weakening,
substitutivity and environment monotonicity for that relation ([S, Sect. 7.3]; [L, Lemma
16]); finally the environment-level relation \ErasesEnv{}, which is [S]'s
dependency-selective \code{erases\_deps} (p. 8:64) rather than the pointwise erasure of
[S, Table 2]. The Lean files are \code{Erasability.lean}, \code{Erases.lean},
\code{ErasesTotal.lean}, \code{Abstract.lean}, \code{ErasesAbstract.lean},
\code{ErasesStrengthen.lean}, \code{ErasesUniform.lean} and \code{ErasesEnv.lean}, all under
\code{LeanToLambdaBox/}; the rule table is mirrored in \code{doc/rules-Erases.md}.

One decision fixes the shape of the transposition: $\Erases\ \mathit{env}\ \mathit{Us}\
\Delta\ e\ t$ is lean4lean's reading relation $\TrExprS\ \mathit{env}\ \mathit{Us}\ \Delta\
e\ e'$ (Chapter~\ref{chap:source}) with the target \code{VExpr} replaced by an \lbox{} term,
the \code{sort} and \code{forallE} arms deleted, and a box rule added. The binder rules then
extend the mixed local context $\Delta$ exactly as \TrExprS{} does, the variable rules carry
\TrExprS{}'s own lookup premises, and the \code{box}, \code{lam} and \code{letE} arms carry
\TrExprS{} witnesses --- so the relation is a bridge into lean4lean's model and inherits its
trust boundary. Two deviations from [S, Fig. 18] are forced by Lean: the locally-nameless
split of \code{erases\_tRel} into a bound-variable and a free-variable rule ([R, Sect. 4.7]),
because the shipping eraser instantiates binders with fresh free variables; and the three
readings of \code{Expr.const}, because Rocq writes \code{tConstruct}, \code{tInd} and
\code{tConst} where Lean writes one node. There is no case, fixpoint or cofixpoint rule:
eliminator applications and top-level recursion are declaration-level phenomena
([R, Sect. 4.1]), handled by \Lower{} in Chapter~\ref{chap:lowering}.

The relation is deliberately non-deterministic --- the box rule competes with every
structural rule and carries no negative side condition, mirroring MetaRocq's
\code{Extract.v} --- so ``the erasure of $e$'' is never well defined here: every statement is
existential, and determinism returns only downstream, and only on first-order answers
(Chapter~\ref{chap:correctness}). Non-vacuity is treated as a standing obligation throughout.
Two items lie outside the main line of argument and are flagged where they occur: the module
\code{ErasesUniform.lean} is outside the capstone's import closure, and \code{doc/coverage.md}
counts its nineteen declarations in the dead-declaration budget. Lemmas whose proofs route
through lean4lean's conversion and unique-typing cluster inherit \code{sorryAx}; each is
marked, and Chapter~\ref{chap:trust} holds the boundary.

\section{Erasability: proofs and type formers}
\label{sec:erases:erasability}

\begin{definition}[Arities]
\label{def:IsArity}
\lean{LeanToLambdaBox.IsArity, LeanToLambdaBox.IsArityUpTo}
\leanok
\uses{def:lean4lean-grounding}
A model type is an \emph{arity} when it is a (possibly empty) $\Pi$-telescope ending in a
sort, $\forall x_1 \dots x_n,\ \mathrm{Sort}\ u$: the inductive \code{IsArity}, with one rule
for \code{sort} and one congruence for \code{forallE}. It is refined to
$\code{IsArityUpTo}\ \mathit{env}\ U\ \Gamma\ A$, ``$A$ is an arity \emph{up to definitional
equality}'': there is an $A'$ with $A \equiv A'$ in $\Gamma$ under $U$ universe parameters and
$\code{IsArity}\ A'$. Only the refined form is invariant under conversion, which is why
\code{Erasable} is stated with it; in Lean it is an existential over a lean4lean conversion
witness. [S] leaves \code{isArity} implicit (p. 8:61) and [L, Def. 1] calls a term with such
a type a \emph{type scheme}.
\end{definition}

\begin{definition}[Erasability]
\label{def:Erasable}
\lean{LeanToLambdaBox.Erasable}
\leanok
\uses{def:IsArity, def:lean4lean-grounding}
$\Erasable\ \mathit{env}\ U\ \Gamma\ e$ holds when the model term $e$ has \emph{some} type $A$
which is either itself a proposition ($A : \mathrm{Sort}\ 0$, so $e$ is a \textbf{proof}) or
an arity up to conversion (so $e$ is a \textbf{type former}). This is [S]'s
$\code{isErasable}\ \Sigma\ \Gamma\ t$ (p. 8:61) transposed to lean4lean's model terms, and
the source-side counterpart of the shipping eraser's test (Chapter~\ref{chap:eraser}). The
existential matters: as in [S], $e$ need only have \emph{some} such type, since unique typing
fails in the source calculus ([L, p. 6]). The Lean definition is phrased with lean4lean's
typing judgement and mentions no \code{Lean.Expr} and no \lbox{} term.
\end{definition}

\begin{lemma}[Syntactic arities are stable]
\label{lem:IsArity-inst}
\lean{LeanToLambdaBox.IsArity.inst, LeanToLambdaBox.IsArity.liftN,
LeanToLambdaBox.IsArity.instL, LeanToLambdaBox.IsArity.of_piBody}
\leanok
\uses{def:IsArity}
Being a syntactic arity is preserved by de Bruijn substitution, by lifting and by
universe-level instantiation; and a telescope whose $\Pi$-body is a sort is an arity, which
is the introduction rule the inductive-declaration metatheory feeds.
\begin{proof}
\leanok
\uses{def:IsArity}
All four by induction on the telescope: each operation fixes \code{.sort} and maps
\code{.forallE} structurally, and \code{piBody} is the telescope's tail.
\end{proof}
\end{lemma}

\begin{remark}
\code{IsArity.instL} is declared in \code{Erases.lean} and \code{IsArity.of\_piBody} in
\code{ErasesTotal.lean} for import-order reasons, away from the rest of the \code{IsArity}
interface; the grouping in this chapter is by concept, not by file. The same holds for
\code{Erasable.mono}, \code{Erasable.instL} and \code{IsArityUpTo.instL} below.
\end{remark}

\begin{lemma}[Arity-up-to-conversion is stable]
\label{lem:IsArityUpTo-inst}
\lean{LeanToLambdaBox.IsArityUpTo.inst, LeanToLambdaBox.IsArityUpTo.weakN,
LeanToLambdaBox.IsArityUpTo.defeqDFC, LeanToLambdaBox.IsArityUpTo.instL}
\leanok
\uses{def:IsArity, def:lean4lean-grounding}
\code{IsArityUpTo} survives instantiating a context entry by a well-typed term, weakening the
context, passing to a definitionally equal context, and instantiating universe levels. Each
Lean statement carries the matching lean4lean context-surgery witness and, for the
substitution form, an ordered environment.
\begin{proof}
\leanok
\uses{lem:IsArity-inst, def:IsArity}
In each case the conversion witness moves by the corresponding upstream conversion lemma and
the syntactic arity by its twin from Lemma~\ref{lem:IsArity-inst}.
\end{proof}
\end{lemma}

\begin{lemma}[Conversion of the type]
\label{lem:IsArityUpTo-defeq}
\lean{LeanToLambdaBox.IsArityUpTo.defeq}
\leanok
\uses{def:IsArity, def:lean4lean-grounding}
If $A''$ is convertible to $A$ and $A$ is an arity up to conversion, so is $A''$. This is the
property the purely syntactic \code{IsArity} lacks and the reason the refined predicate
exists; the Lean statement asks for a well-formed environment and context.
\begin{proof}
\leanok
\uses{def:IsArity, asm:lean4lean-trust}
Compose the two conversions by transitivity and keep the syntactic witness. Transitivity lies
in lean4lean's unique-typing cluster, so the measured axiom footprint contains
\code{sorryAx}.
\end{proof}
\end{lemma}

\begin{lemma}[Irrelevance survives reduction]
\label{lem:Erasable-defeq}
\lean{LeanToLambdaBox.Erasable.defeq}
\leanok
\uses{def:Erasable, def:lean4lean-grounding}
\textbf{The box-soundness core.} If $e$ is erasable and $e \equiv e'$, then $e'$ is erasable.
A reduction step is a definitional equality, so an irrelevant term stays irrelevant as it
reduces, which is what makes erasing it to $\Box$ sound in the simulation of
Chapter~\ref{chap:correctness}. The Lean statement asks for a well-formed environment and
context. Its informal ancestor is [L, Lemma 2]: \code{Prop}-sortedness and being a type
scheme are preserved by reduction.
\begin{proof}
\leanok
\uses{def:Erasable, asm:lean4lean-trust}
Keep the same type $A$ and move the typing witness along the conversion; the proof/arity
disjunct concerns $A$ alone and carries over unchanged. The measured axiom footprint contains
\code{sorryAx}, inherited through the conversion cluster.
\end{proof}
\end{lemma}

\begin{lemma}[Stability of irrelevance]
\label{lem:Erasable-weakN}
\lean{LeanToLambdaBox.Erasable.weakN, LeanToLambdaBox.Erasable.inst,
LeanToLambdaBox.Erasable.defeqDFC, LeanToLambdaBox.Erasable.instL,
LeanToLambdaBox.Erasable.mono}
\leanok
\uses{def:Erasable, def:lean4lean-grounding}
Irrelevance is preserved by weakening the context, by substituting a well-typed term, by
passing to a definitionally equal context, by universe instantiation, and by extending the
environment. These are exactly the five moves the transport results of
Sections~\ref{sec:erases:transport} and~\ref{sec:erases:strengthen} perform, one per
\code{box} arm. The Lean statements carry the matching surgery witnesses; the universe form
additionally asks that the substituted levels be well formed in the target scope.
\begin{proof}
\leanok
\uses{lem:IsArityUpTo-inst, def:Erasable}
Move the typing witness by the corresponding upstream lemma and the arity disjunct by its
twin from Lemma~\ref{lem:IsArityUpTo-inst}. The sort $\mathrm{Sort}\ 0$ is fixed by lifting,
substitution and level instantiation, so the proof disjunct carries over verbatim.
\end{proof}
\end{lemma}

\begin{theorem}[Box propagation through application]
\label{thm:Erasable-app}
\lean{LeanToLambdaBox.Erasable.app}
\leanok
\uses{def:Erasable, def:lean4lean-grounding}
If the head $f$ is irrelevant and the application $f\,a$ is well typed, then $f\,a$ is
irrelevant. This is the type-level content of MetaRocq's \code{eval\_box} rule, that is, of
[S, Sect. 7.1] and [L, Def. 5]: $\Box\ u \to \Box$. If $f$ is a proof, its type is convertible
to $\forall x : A,\ B$ living in $\mathrm{Sort}\ 0$, so $\mathrm{imax}\ u\ v \approx 0$ forces
$v \approx 0$ and $f\,a : B[a]$ is again a proof; if $f$ is a type former, $B$ is an arity up
to conversion, hence so is $B[a]$. The Lean statement takes the function type and the
argument's typing as separate premises. It is the source-side \emph{licence} for the \lbox{}
box rule, which itself lives in Chapter~\ref{chap:lambdabox}.
\begin{proof}
\leanok
\uses{def:Erasable, def:IsArity, lem:IsArityUpTo-defeq, lem:IsArityUpTo-inst,
asm:lean4lean-trust}
Retype $f$ against its function type by uniqueness of typing, then split on the
\code{Erasable} disjunct. In the proof case invert $\mathrm{imax}\ u\ v \approx 0$ by the
sort-inversion lemmas to get $B : \mathrm{Sort}\ 0$; in the type-former case push the arity
through the substitution by Lemma~\ref{lem:IsArityUpTo-inst} and through the conversion by
Lemma~\ref{lem:IsArityUpTo-defeq}. Uniqueness and sort inversion are lean4lean \code{sorry}
sites; the measured axiom footprint contains \code{sorryAx}.
\end{proof}
\end{theorem}

\begin{definition}[Relevance of an inductive type former]
\label{def:InformativeInd}
\lean{LeanToLambdaBox.InformativeInd, LeanToLambdaBox.vResultSort,
LeanToLambdaBox.neverZeroB}
\leanok
\uses{def:lean4lean-grounding}
$\code{InformativeInd}\ \mathit{env}\ I$ holds when \emph{env} declares $I$ and the result
sort of $I$'s declared type never evaluates to $\mathrm{Prop}$ under any instantiation of its
universe parameters. The auxiliary \code{vResultSort} reads the result level off a
$\Pi$-telescope; \code{neverZeroB} decides ``this level is never zero'' (a successor never is,
a parameter may be instantiated by zero, a \code{max} needs one non-zero side, an \code{imax}
is non-zero exactly when its right argument is). The criterion is \textbf{semantic} --- it is
lean4lean's own never-zero predicate, the one its kernel theory reads for large elimination
--- not the syntactic ``the result sort is a successor''. This is the one premise
[S, Fig. 18] does not spell (design tag N18); it gates the \code{proj} rule of
Definition~\ref{def:Erases} and the source $\iota$ rule of Chapter~\ref{chap:source}, and
\code{neverZeroB} is what makes it decidable for the fragment checker.
\end{definition}

\begin{lemma}[The relevance test is sound and complete]
\label{lem:neverZeroB_sound}
\lean{LeanToLambdaBox.neverZeroB_sound, LeanToLambdaBox.neverZeroB_complete}
\leanok
\uses{def:InformativeInd}
$\code{neverZeroB}\ l = \mathit{true}$ implies that $l$ never evaluates to zero, and
$\code{neverZeroB}\ l = \mathit{false}$ implies that the all-zero instantiation witnesses that
it does.
\begin{proof}
\leanok
\uses{def:InformativeInd}
Induction on the level, unfolding level evaluation at each constructor; the \code{imax} case
uses $\mathrm{imax}\ u\ 0 = 0$.
\end{proof}
\end{lemma}

\begin{remark}
\code{neverZeroB\_sound} is one of the few declarations of this chapter whose footprint the
committed fixture \code{test/ledger.expected} records: \code{propext}, \code{Classical.choice}
and \code{Quot.sound}, so no \code{sorryAx}, and hence no trust edge on this node.
\end{remark}

\begin{lemma}[Relevance is monotone]
\label{lem:InformativeInd-mono}
\lean{LeanToLambdaBox.InformativeInd.mono, LeanToLambdaBox.informativeInd_of_succ}
\leanok
\uses{def:InformativeInd}
Relevance survives environment extension, which is what the \code{proj} arm of
Theorem~\ref{thm:Erases-mono} needs; and every producer of the syntactic successor criterion
produces the semantic one, which lets the first-order layer of
Chapter~\ref{chap:correctness} keep its successor clause and still supply relevance.
\begin{proof}
\leanok
\uses{def:InformativeInd}
Transport the constant lookup along the extension; a successor level evaluates to a
successor, hence never to zero.
\end{proof}
\end{lemma}

\begin{lemma}[The two criteria are different]
\label{lem:informativeInd_prod}
\lean{LeanToLambdaBox.informativeInd_prod, LeanToLambdaBox.not_informativeInd_and,
LeanToLambdaBox.prodVTy, LeanToLambdaBox.andVTy}
\leanok
\uses{def:InformativeInd}
On an environment declaring \code{Prod} at the translated type $\mathrm{Sort}\ u \to
\mathrm{Sort}\ v \to \mathrm{Sort}\ (\max(u{+}1, v{+}1))$ the semantic criterion accepts
\code{Prod}, which the syntactic successor test rejects; on an environment declaring
\code{And} at $\mathrm{Prop} \to \mathrm{Prop} \to \mathrm{Prop}$ it rejects \code{And}. Each
Lean statement is conditional on the environment declaring the stated literal type, which is
all a hand-built model environment can offer. They are the separating witnesses that justify
the semantic criterion.
\begin{proof}
\leanok
\uses{lem:neverZeroB_sound, def:InformativeInd}
Compute the result sort of the two literal types; for \code{Prod} discharge never-zeroness
through Lemma~\ref{lem:neverZeroB_sound}, and for \code{And} the result level is zero.
\end{proof}
\end{lemma}

\section{Reading the kernel environment: blocks, constructors, constants}
\label{sec:erases:readings}

\begin{definition}[Block data in \texorpdfstring{\lbox{}}{lambda-box} coordinates]
\label{def:IndInfo}
\lean{LeanToLambdaBox.IndInfo, LeanToLambdaBox.IndInfo.block,
LeanToLambdaBox.indBlockKername, LeanToLambdaBox.ctorFieldCounts}
\leanok
\uses{def:InductiveId, def:Kername, def:lean4lean-grounding}
$\code{IndInfo}\ \mathit{env}\ I\ \mathit{iid}\ \mathit{np}\ \mathit{nfs}$ is \emph{env}'s
data for the inductive $I$ \textbf{in \lbox{} coordinates}: some declaration list, well formed
for an environment $\mathit{env}_0 \le \mathit{env}$, declares a mutual block with
$\mathit{np}$ parameters whose $\mathit{iid}$-indexed type former is $I$; $\mathit{iid}$'s
block kername is exactly \code{indBlockKername} of that block's names, the kername the eraser
mints by concatenating the type formers' names and reading the result as a root kername; and
$\mathit{nfs}$ lists $I$'s constructors' field counts, \code{ctorFieldCounts} being the
$\Pi$-binders beyond the parameters. A model environment stores no declarations, so a block
must be read off a well-formed declaration list; and the list is required to be well formed
for an environment \emph{below} \emph{env} rather than for \emph{env} itself, which is
precisely what makes the predicate monotone. Formally it is a one-field structure, the field
being that existential. \code{indBlockKername} is the only place the frontend's naming
convention enters the specification.
\end{definition}

\begin{definition}[Block data in source coordinates]
\label{def:IndArity}
\lean{LeanToLambdaBox.IndArity}
\leanok
\uses{def:IndInfo, def:lean4lean-grounding}
$\code{IndArity}\ \mathit{env}\ I\ \mathit{np}\ \mathit{nfs}$ is the same content as
Definition~\ref{def:IndInfo} with the \lbox{} identifier dropped: the parameter count and the
per-constructor field counts, read off the same kind of declaration list below \emph{env}. It
is the form a rule about \code{Lean.Expr} evaluation may read, since it mentions no target
syntax.
\end{definition}

\begin{definition}[Declared in the environment's own list]
\label{def:IndDeclOf}
\lean{LeanToLambdaBox.IndDeclOf}
\leanok
\uses{def:lean4lean-grounding}
$\code{IndDeclOf}\ \mathit{env}\ I$ says that $I$ is declared by a block of
\textbf{\emph{env}'s own} declaration list --- one extension step stronger than
\code{IndInfo}, which exhibits a block below \emph{env}. It is deliberately \emph{not}
monotone in \emph{env}, and therefore deliberately on no rule of \Erases{}; it is the
\code{declared\_inductive} $\Sigma$ conjunct of [S]'s \code{erases\_deps}, supplied instead by
\ErasesEnv{}'s block clause (Definition~\ref{def:ErasesEnv}).
\end{definition}

\begin{definition}[The constructor reading]
\label{def:CtorOf}
\lean{LeanToLambdaBox.CtorOf}
\leanok
\uses{def:lean4lean-grounding}
$\code{CtorOf}\ \mathit{env}\ c\ I\ k$ says that $c$ is the $k$-th constructor of the
inductive type $I$, read off a well-formed declaration list below \emph{env} --- the same
reading \code{IndInfo} takes. This is the premise that turns a \code{Lean.Expr.const} node
into Rocq's \code{tConstruct}.
\end{definition}

\begin{definition}[The plain-constant reading]
\label{def:ConstOrigin}
\lean{LeanToLambdaBox.ConstOrigin, LeanToLambdaBox.VDeclDefines}
\leanok
\uses{def:lean4lean-grounding}
$\code{ConstOrigin}\ \mathit{env}\ c$ is Rocq's \code{tConst} reading: $c$ is introduced as a
definition, an opaque constant, an example, a member of a mutual definition block, or an
axiom, with the declaration exhibited off a declaration list below \emph{env}.
\code{VDeclDefines} is the per-declaration test; a block introduces type names and
constructors, and the quotient declaration is outside the fragment, so neither counts. The
predicate is \textbf{positive}: it exhibits the declaration rather than excluding the other
two readings, so an introduction site discharges it from the declaration list it already has.
The exclusion direction --- that this reading rules the other two out --- is a theorem of the
kernel theory and enters the development as an upstream ask (Chapter~\ref{chap:trust}), not as
a premise here.
\end{definition}

\begin{lemma}[The readings are monotone]
\label{lem:IndInfo-mono}
\lean{LeanToLambdaBox.IndInfo.mono, LeanToLambdaBox.IndArity.mono,
LeanToLambdaBox.CtorOf.mono, LeanToLambdaBox.ConstOrigin.mono,
LeanToLambdaBox.IndInfo.of_wf', LeanToLambdaBox.ConstOrigin.of_axiom}
\leanok
\uses{def:IndInfo, def:IndArity, def:CtorOf, def:ConstOrigin}
Each of the four readings survives an extension $\mathit{env} \le \mathit{env}'$. Two
introduction rules are recorded alongside: \code{IndInfo.of\_wf'} builds block data from an
explicit declaration list, and \code{ConstOrigin.of\_axiom} says that a constant an
environment declares at the type of a closed axiom has that axiom's declaration as its origin
--- the cheapest introduction, and the one a hand-built environment uses.
\begin{proof}
\leanok
\uses{def:IndInfo, def:CtorOf, def:ConstOrigin}
Each witness already carries the bound $\mathit{env}_0 \le \mathit{env}$; compose it with
$\mathit{env} \le \mathit{env}'$ by transitivity and keep everything else.
\end{proof}
\end{lemma}

\begin{remark}
This is the sole reason the bound $\mathit{env}_0 \le \mathit{env}$ appears in all four
predicates, and the reason \code{IndDeclOf} (Definition~\ref{def:IndDeclOf}) is on no rule.
\end{remark}

\begin{lemma}[Exchanging block data between coordinates]
\label{lem:IndInfo-arity}
\lean{LeanToLambdaBox.IndInfo.arity, LeanToLambdaBox.IndArity.indInfo,
LeanToLambdaBox.CtorOf.indInfo, LeanToLambdaBox.CtorOf.indArity}
\leanok
\uses{def:IndInfo, def:IndArity, def:CtorOf}
The \lbox{} identifier is the only thing \code{IndInfo} has beyond \code{IndArity}, and
conversely source-coordinate block data \emph{exhibits} an identifier, because it names a
position in its own block: the identifier is not extra data. Moreover a constructor's own
declaration is its type's block, so all the block data the \code{ctor} rule needs is derivable
from \code{CtorOf} alone. These are what let source-side rules and target-side rules exchange
block data.
\begin{proof}
\leanok
\uses{def:IndInfo, def:CtorOf}
Forgetting the index equation is list membership; recovering an index from membership is list
indexing, after which the block identifier is \code{indBlockKername} of the block's names and
its defining equation is definitional. The constructor forms do the same starting from the
constructor's own block.
\end{proof}
\end{lemma}

\section{The erasure relation and its eleven rules}
\label{sec:erases:relation}

\begin{definition}[The erasure relation]
\label{def:Erases}
\lean{LeanToLambdaBox.Erases, LeanToLambdaBox.Erases.box, LeanToLambdaBox.Erases.bvar,
LeanToLambdaBox.Erases.fvar, LeanToLambdaBox.Erases.ctor, LeanToLambdaBox.Erases.const,
LeanToLambdaBox.Erases.app, LeanToLambdaBox.Erases.lam, LeanToLambdaBox.Erases.letE,
LeanToLambdaBox.Erases.proj, LeanToLambdaBox.Erases.lit, LeanToLambdaBox.Erases.mdata}
\leanok
\uses{def:Erasable, def:CtorOf, def:IndInfo, def:ConstOrigin, def:InformativeInd,
def:LBTerm, def:Kername, def:lean4lean-grounding}
$\Erases\ \mathit{env}\ \mathit{Us}\ \Delta\ e\ t$ relates a source term
$e : \code{Lean.Expr}$, read in the mixed local context $\Delta$ exactly as \TrExprS{} reads
it, to a \lbox{} image $t$. There are eleven rules:

\medskip
\noindent\begin{tabular}{lllp{3.7cm}}
rule & source & image & premises \\
\hline
\code{box} & any $e$ & $\Box$ & $e$ translates to $\mathit{ve}$, and $\mathit{ve}$ is
\code{Erasable} \\
\code{bvar} & $\code{bvar}\ i$ & $\code{bvar}\ i$ & $\Delta$ binds the index $i$ \\
\code{fvar} & $\code{fvar}\ x$ & $\code{fvar}\ x$ & $\Delta$ binds the variable $x$ \\
\code{ctor} & $\code{const}\ c\ \mathit{us}$ & $\code{construct}\ \mathit{iid}\ k\ []$ &
$\code{CtorOf}\ \mathit{env}\ c\ I\ k$, and \code{IndInfo} for $I$ \\
\code{const} & $\code{const}\ c\ \mathit{us}$ & $\code{const}\ (\code{toKername}\ c)$ &
\emph{env} declares $c$, and $\code{ConstOrigin}\ \mathit{env}\ c$ \\
\code{app} & $f\ a$ & $f'\ a'$ & both components erase \\
\code{lam} & $\lambda x{:}A.\,b$ & $\lambda x.\,b'$ & $A$ translates; $b$ erases in $\Delta$
extended by a $\lambda$ entry \\
\code{letE} & $\code{let}\ x{:}A := v;\ b$ & $\code{let}\ x := v';\ b'$ & $A$ and $v$
translate, $v$ erases, $b$ erases under a \code{let} entry \\
\code{proj} & $\code{proj}\ S\ i\ e$ & $\code{proj}\ \langle \mathit{iid},\mathit{np},i
\rangle\ t$ & $\code{IndInfo}\ \mathit{env}\ S\ \mathit{iid}\ \mathit{np}\ [\mathit{nf}]$,
$i < \mathit{nf}$, and $\code{InformativeInd}\ \mathit{env}\ S$ \\
\code{lit} & $\code{lit}\ l$ & $t$ & \emph{env} contains the literal's unfolding
constants (\code{ContainsLits}), and $l$'s one-step kernel unfolding erases to $t$ \\
\code{mdata} & $\code{mdata}\ d\ e$ & $t$ & $e$ erases to $t$ \\
\end{tabular}
\medskip

Three points deserve emphasis. The \code{ctor} node carries a \textbf{literally empty}
argument list: it sits at the bare head, and arguments arrive through \code{app}. Where Rocq
has three term nodes --- and [S, Fig. 18] therefore three rules --- Lean has one,
\code{Expr.const}, so \code{ctor} and \code{const} carry the premise saying which reading
applies; an inductive type \emph{name} gets no rule at all, because it is erasable and $\Box$
is its only image (Lemma~\ref{lem:Erases-indInfo_erasable}). And no rule produces a case or a
fixpoint node. The Lean statement additionally records that the two binder rules keep the
\textbf{source} binder name in the image, and that \code{Expr.letE}'s non-dependence flag is
ignored.
\end{definition}

\begin{remark}
The arm list is pinned by a compile-time \code{\#guard} against an anonymous elaborator that
reads the relation's constructor names out of the environment, and \code{lake exe hygiene
--tables} checks that every arm appears in \code{doc/rules-Erases.md}; adding or removing a
rule breaks the build. The two \code{Expr.const} rules are \emph{not} symmetric: \code{const}
additionally carries a constant-lookup premise, while \code{ctor} carries none, because
\code{CtorOf} and \code{IndInfo} already exhibit the declaring block.
\end{remark}

\begin{remark}[Correspondence with {[S, Fig. 18]}]
\code{erases\_box} becomes \code{box}; \code{erases\_tRel} splits into \code{bvar} and
\code{fvar}; \code{erases\_tLambda}, \code{erases\_tLetIn}, \code{erases\_tApp},
\code{erases\_tConst}, \code{erases\_tConstruct} and \code{erases\_tProj} become \code{lam},
\code{letE}, \code{app}, \code{const}, \code{ctor} and \code{proj}. \code{erases\_tCase} and
the subsingleton-elimination rule have no counterpart, Lean having no case node: they are
\code{Lower.elimApp}'s business. \code{erases\_tFix} has none, top-level recursion being
\code{Lower.fixConst} and \code{Lower.fixBody}; \code{erases\_tCoFix} none, Lean having no
coinduction. An inductive type name gets no rule. In the other direction \code{mdata} and
\code{lit} are added, and [S]'s sort and $\forall$ premises are absorbed into \code{box}
(Lemma~\ref{lem:Erases-sort_erasable}). [L]'s counterpart is not the erasure function but the
invariant of [L, Def. 10]: the \lbox{} term differs from the source only by boxing, and every
boxed position is erasable.
\end{remark}

\begin{definition}[The box alternative]
\label{def:ErasesBox}
\lean{LeanToLambdaBox.ErasesBox}
\leanok
\uses{def:Erasable, def:LBTerm, def:lean4lean-grounding}
$\code{ErasesBox}\ \mathit{env}\ \mathit{Us}\ \Delta\ e\ t$ abbreviates the alternative every
inversion lemma shares: $e$ has a model translation that is \code{Erasable}, and $t$ is
$\Box$.
\end{definition}

\begin{lemma}[The inversion kit]
\label{lem:Erases-sort_inv}
\lean{LeanToLambdaBox.Erases.sort_inv, LeanToLambdaBox.Erases.forallE_inv,
LeanToLambdaBox.Erases.bvar_inv, LeanToLambdaBox.Erases.fvar_inv,
LeanToLambdaBox.Erases.const_inv, LeanToLambdaBox.Erases.ctor_inv,
LeanToLambdaBox.Erases.app_inv, LeanToLambdaBox.Erases.lam_inv,
LeanToLambdaBox.Erases.letE_inv, LeanToLambdaBox.Erases.proj_inv,
LeanToLambdaBox.Erases.lit_inv, LeanToLambdaBox.Erases.mdata_inv}
\leanok
\uses{def:Erases, def:ErasesBox, def:CtorOf, def:IndInfo, def:ConstOrigin,
def:InformativeInd}
One lemma per source head, each a disjunction whose first alternative is \code{ErasesBox}. A
sort and a $\Pi$-type admit \emph{only} that alternative. A bound or free variable admits,
besides box, the lookup premise with the same variable as image. A constant admits
\textbf{three} alternatives, one per reading: box; the \code{ctor} image with \code{CtorOf}
and \code{IndInfo} exhibited; the \code{const} image with the constant lookup and
\code{ConstOrigin}. Nothing here excludes two of the three --- that is the kernel theory's
business --- so a consumer that knows its reading supplies the corresponding premise and
discards the others. One lemma is keyed on the \textbf{target}: if any source erases to
$\code{construct}\ \mathit{iid}\ k\ \mathit{args}$ then $\mathit{args} = []$ and the block
data is exhibited, the source not being pinned because \code{lit} and \code{mdata} pass an
image through. For the same reason the literal lemma does \emph{not} pin the image.
\begin{proof}
\leanok
\uses{def:Erases, def:ErasesBox}
Each is a case analysis on the relation, except the target-keyed one, which generalises the
image and inducts so that the \code{lit} and \code{mdata} arms can pass it through.
\end{proof}
\end{lemma}

\begin{remark}
This is the elimination interface every downstream proof uses instead of raw case analysis.
The sort and $\Pi$ lemmas are also what \code{doc/panics.md} consumes to argue that the
shipping eraser's \code{unreachable!} at those heads is unreachable
(Chapter~\ref{chap:eraser}).
\end{remark}

\begin{theorem}[Both constant readings fire]
\label{thm:erases_ctor_fires}
\lean{LeanToLambdaBox.erases_ctor_fires, LeanToLambdaBox.erases_const_fires,
LeanToLambdaBox.blkDecl_wf, LeanToLambdaBox.blk_wf', LeanToLambdaBox.blk_ctorOf,
LeanToLambdaBox.blk_indInfo, LeanToLambdaBox.blk_constOrigin}
\leanok
\uses{def:Erases}
On a hand-built one-inductive, one-definition kernel environment --- a \code{Prop}-valued type
former with one nullary constructor, its eliminator, and a definition bound to that
constructor --- the constructor reading and the constant reading each actually fire: the
fixture's constructor erases to its \lbox{} constructor node and the fixture's definition to
its kername, with every premise discharged from the fixture's own declaration list. The Lean
statements are universally quantified over the universe scope, the local context and the
constant's level arguments.
\begin{proof}
\leanok
\uses{def:Erases, def:CtorOf, def:IndInfo, def:ConstOrigin, lem:IndInfo-mono}
Discharge the full inductive-declaration well-formedness of the fixture's block --- some
twenty obligations, including a typed $\iota$ rule --- by explicit typing derivations and
decision procedures; stage the environment through the type, constructor and recursor
additions with definitional equations to reach well-formedness of the two-declaration list;
then read \code{CtorOf}, \code{IndInfo} and \code{ConstOrigin} off that list and apply the two
rules.
\end{proof}
\end{theorem}

\begin{remark}
This is \emph{unconditional} non-vacuity --- no hypothesis about the source --- which is worth
stating, because the corresponding witness for \ErasesEnv{}
(Theorem~\ref{thm:demoEnv_erasesEnv}) is conditional. The fixture's type former is
\code{Prop}-valued, so the \code{proj} rule provably does \emph{not} fire at it, which is the
point of the next theorem.
\end{remark}

\begin{theorem}[A boxed discriminant is stuck]
\label{thm:erases_proj_needs_informative}
\lean{LeanToLambdaBox.erases_proj_needs_informative}
\leanok
\uses{def:WcbvEval, def:eraseFlags, def:GlobalDeclarations, def:InductiveId}
In \textbf{any} \lbox{} environment and at \textbf{any} projection, the term
$\code{proj}\ p\ \Box$ has no value under the erasure flag configuration \code{eraseFlags}:
the applied-form projection rule demands a constructor spine, the block-form rule is off, and
the propositional rule needs a flag the erasure-correctness statement does not enable.
Consequently the relevance premise on the \code{proj} rule is not a convenience: without it
the rule would fire at a projection out of a \code{Prop}-valued structure, whose discriminant
is a proof and may erase to $\Box$, and the corresponding arm of the simulation would be
\textbf{refutable}, not merely unprovable.
\begin{proof}
\leanok
\uses{def:InformativeInd, def:WcbvEval}
A $\Box$ evaluates only to itself, so its spine head is $\Box$ and not a constructor, which
kills the applied-form rule; the two remaining projection rules are refuted by deciding the
flag record, whose block and propositional flags are off.
\end{proof}
\end{theorem}

\begin{remark}
This is the chapter's model for how an added premise is justified: by a refutation rather than
by assertion. [S, Fig. 18] needs no such premise because Rocq's typing forbids a projection
out of \code{Prop}, whereas Lean's \code{Expr.proj} occurs there --- \code{h.1} on
\code{And}. The rule table argues that the premise removes no program, since every field of a
propositional structure is a proof and \code{box} is its rule.
\end{remark}

\section{Monotonicity and totality}
\label{sec:erases:total}

\begin{theorem}[Environment weakening]
\label{thm:Erases-mono}
\lean{LeanToLambdaBox.Erases.mono}
\leanok
\uses{def:Erases}
An erasure derivation survives an extension $\mathit{env} \le \mathit{env}'$ of the model
environment, with source and target unchanged. This is [S, Sect. 7.3]'s global weakening for
the erasure relation, and the reason all four environment readings carry the bound
$\mathit{env}_0 \le \mathit{env}$.
\begin{proof}
\leanok
\uses{def:Erases, lem:Erasable-weakN, lem:IndInfo-mono, lem:InformativeInd-mono}
Induction on the derivation: each premise is monotone in its own right. Translations move by
lean4lean's own monotonicity, \code{Erasable.mono} handles \code{box}, \code{CtorOf.mono} and
\code{IndInfo.mono} handle \code{ctor}, the constant lookup with \code{ConstOrigin.mono}
handles \code{const}, \code{IndInfo.mono} with \code{InformativeInd.mono} handles \code{proj},
and the literal premise moves by its own monotonicity lemma.
\end{proof}
\end{theorem}

\begin{lemma}[Sorts and $\Pi$-types are erasable]
\label{lem:Erases-sort_erasable}
\lean{LeanToLambdaBox.Erases.sort_erasable, LeanToLambdaBox.Erases.forallE_erasable}
\leanok
\uses{def:Erasable, def:lean4lean-grounding}
A sort is a type former --- its type is the next sort up, an arity --- and so is a $\Pi$-type,
whose type is the \code{imax} of its parts' sorts. Hence the box rule covers both heads and no
structural rule for them is missing: this is how [S, Fig. 18]'s sort and $\forall$ premises
are absorbed. The $\Pi$ statement asks for a well-formed environment and context; the sort
statement binds an environment premise that its proof never uses, as its doc-comment says.
\begin{proof}
\leanok
\uses{def:Erasable, def:IsArity}
Invert the translation: the model term's type is a sort, respectively an \code{imax} sort,
which is a one-node arity, so the arity disjunct holds by reflexivity of conversion. The $\Pi$
case needs the context's well-formedness only to supply the levels' well-formedness.
\end{proof}
\end{lemma}

\begin{lemma}[A type former's declared type is an arity]
\label{lem:IndInfo-constant_isArity}
\lean{LeanToLambdaBox.IndInfo.constant_isArity, LeanToLambdaBox.wf'_induct_origin}
\leanok
\uses{def:IndInfo, def:IsArity}
A block occurring in a well-formed declaration list was added by a well-formed declaration
step, and the environment that step produced sits below the list's own; from which the
declared type of an inductive type former is an arity. This is the cost of having no
\code{tInd} rule: it is what makes an inductive type name erasable.
\begin{proof}
\leanok
\uses{def:IndInfo, lem:IsArity-inst}
Peel the declaration list back to the step that added the block. That step binds the type
former's name to exactly the type the declaration's universe condition forces to end in a
sort; conclude with \code{IsArity.of\_piBody}.
\end{proof}
\end{lemma}

\begin{lemma}[An inductive type name is erasable]
\label{lem:Erases-indInfo_erasable}
\lean{LeanToLambdaBox.Erases.indInfo_erasable}
\leanok
\uses{def:Erasable, def:IndInfo, def:lean4lean-grounding}
A constant that \code{IndInfo} reads as an inductive type former is erasable, because its
declared type is an arity and the arity survives the universe instantiation the translation
performs. So $\Box$ is its only image, and the third reading of \code{Expr.const} is closed
without a rule. The Lean statement asks for a well-formed environment and context and takes
the constant's translation as a premise.
\begin{proof}
\leanok
\uses{lem:IndInfo-constant_isArity, lem:IsArity-inst, def:Erasable}
Invert the translation of the constant, identify its type with the declared one, and push the
arity through level instantiation.
\end{proof}
\end{lemma}

\begin{remark}
This is also the concrete reason the \code{const} rule cannot be premise-free: a ten-rule
relation without \code{ConstOrigin} would send \code{Nat} to its kername, which the eraser
emits only as an inductive declaration, and on which the target's $\delta$ rule is stuck.
\end{remark}

\begin{definition}[Block data at every projection]
\label{def:ProjInfo}
\lean{LeanToLambdaBox.ProjInfo, LeanToLambdaBox.ProjInfo.toConstructor}
\leanok
\uses{def:IndInfo, def:InformativeInd}
$\code{ProjInfo}\ \mathit{env}\ e$ says that at every projection head of $e$ the projected
structure carries block data --- a single-constructor \code{IndInfo} with the field index in
range --- and is relevant. All other heads hold trivially, and the congruences are the
computational ones: binder \emph{types} are not traversed, since they reach erasure only as
\TrExprS{} witnesses. A literal needs no premise, because its one-step kernel unfolding is
projection-free. The predicate carries exactly what the \code{proj} rule demands and what
lean4lean's own projection premise does \emph{not} yield; its purpose is to carve the
projection case out of the totality theorem without importing lean4lean's unproven projection
cluster for one case. Downstream it is discharged by the fragment checker of
Chapter~\ref{chap:source}.
\end{definition}

\begin{theorem}[Totality]
\label{thm:Erases-exists_of_trExprS_of_projInfo}
\lean{LeanToLambdaBox.Erases.exists_of_trExprS_of_projInfo}
\leanok
\uses{def:Erases, def:ProjInfo, def:CtorOf, def:IndInfo, def:ConstOrigin,
def:lean4lean-grounding}
Every source term that lean4lean can read has at least one \lbox{} image, under a well-formed
environment and local context, provided its projection heads carry block data and relevance
and its declared constants are classified: the Lean statement binds a hypothesis saying that
every constant the environment declares is a constructor, or an inductive type former, or a
plain constant. Every head takes its structural rule, except a sort, a $\Pi$-type and an
inductive type name, where the box rule fires.
\begin{proof}
\leanok
\uses{lem:Erases-sort_erasable, lem:Erases-indInfo_erasable, lem:IndInfo-arity,
def:ProjInfo, def:Erases}
Induction on the translation. At a sort and a $\Pi$-type the box rule fires by
Lemma~\ref{lem:Erases-sort_erasable}. At a constant, split with the classification
hypothesis: a constructor takes the \code{ctor} rule, its block data supplied by
Lemma~\ref{lem:IndInfo-arity}; a type former takes the box rule by
Lemma~\ref{lem:Erases-indInfo_erasable}; a plain constant takes the \code{const} rule. The
binder cases extend the context exactly as the translation does, the literal case uses that a
literal's unfolding is projection-free, and the projection case reads its two premises off
\code{ProjInfo}.
\end{proof}
\end{theorem}

\begin{remark}
This is the ``erasure is defined everywhere on the fragment'' half of the specification, the
analogue of [S]'s \code{erases\_erase} (p. 8:62) at the level of existence. It is \emph{not}
the statement that the shipping function's graph lies in the relation, which is the bridge of
Chapter~\ref{chap:refinement}. The classification premise is a kernel-theory fact; it is a
hypothesis of this theorem alone, discharged elsewhere from the upstream-ask bundle
(Chapter~\ref{chap:trust}), and never a premise of \Erases{} itself.
\end{remark}

\section{Closing a free variable: the \texorpdfstring{\lbox{}}{lambda-box}-side algebra}
\label{sec:erases:abstract}

\begin{definition}[Occurrence of a free variable in a \texorpdfstring{\lbox{}}{lambda-box} term]
\label{def:hasFVar}
\lean{LeanToLambdaBox.hasFVar, LeanToLambdaBox.hasFVarArgs, LeanToLambdaBox.hasFVarAlts,
LeanToLambdaBox.hasFVarDefs, LeanToLambdaBox.hasFVarArgs_iff,
LeanToLambdaBox.hasFVarAlts_iff, LeanToLambdaBox.hasFVarDefs_iff}
\leanok
\uses{def:LBTerm}
$\code{hasFVar}\ x\ t$ is structural occurrence of the free variable $x$ in the \lbox{} term
$t$. The recursion mirrors the closing operator \code{toBvar}'s exactly, with dedicated
mutually recursive helpers for the three nested lists --- constructor arguments, case
alternatives, fixpoint definitions --- because the naive existential clause is rejected by the
termination checker; three \code{\_iff} lemmas recover the expected existential form. The
predicate is consumed well beyond this chapter.
\end{definition}

\begin{lemma}[Closing: the occurrence laws]
\label{lem:toBvar_eq_of_not_hasFVar}
\lean{LeanToLambdaBox.toBvar_eq_of_not_hasFVar, LeanToLambdaBox.abstract_eq_of_not_hasFVar,
LeanToLambdaBox.hasFVar_toBvar_of, LeanToLambdaBox.fvarId_beq_iff_eq}
\leanok
\uses{def:hasFVar, def:toBvar}
Closing over a variable that does not occur is the identity; and closing never \emph{creates}
an occurrence --- whatever occurs in $\code{toBvar}\ y\ \mathit{lvl}\ t$ occurred in $t$
already and is not $y$, which has been replaced by a de Bruijn index everywhere. A small
reflection lemma for the derived equality test on free-variable identifiers bridges the
Boolean test in \code{toBvar} and the propositional clause in \code{hasFVar}; it is proved
locally because Lean core ships no lawfulness instance for that test.
\begin{proof}
\leanok
\uses{def:hasFVar}
Structural induction, mutual with the three list traversals, using the reflection lemma at the
free-variable leaf.
\end{proof}
\end{lemma}

\begin{remark}
The module documentation records that these lemmas were unprovable while \code{toBvar} was a
\code{partial def}; de-partialising it into a structural definition is what made them
available. In the bridge of Chapter~\ref{chap:refinement} the no-op lemma is what makes the
shipping eraser's abstraction step vanish on subterms that never mention the freshly opened
variable.
\end{remark}

\begin{lemma}[Closing commutes with shifting and with itself]
\label{lem:toBvar_shift}
\lean{LeanToLambdaBox.toBvar_shift, LeanToLambdaBox.toBvar_toBvar,
LeanToLambdaBox.toBvarDefs_length, LeanToLambdaBox.shiftDefs_length,
LeanToLambdaBox.abstract_eq, LeanToLambdaBox.toBvarArgs_eq_map,
LeanToLambdaBox.toBvarAlts_eq_map, LeanToLambdaBox.toBvarDefs_eq_map}
\leanok
\uses{def:toBvar, def:LBTerm-shift}
Closing and de Bruijn shifting commute, provided the insertion level is at or above the
shift's cutoff --- the configuration the bridge meets, where abstraction happens at the level
of the opened binder and shifts happen below it --- with the insertion level moving by the
shift amount on the shifted side. Two closings at distinct variables commute with \emph{no}
index adjustment at all, because closing never rewrites existing indices; that is what the
bridge cases which open several free variables at once and close them in sequence need. The
three list traversals of the closing operator are ordinary maps, closing and shifting preserve
the length of a fixpoint block, and \code{abstract} is closing at level $0$.
\begin{proof}
\leanok
\uses{lem:toBvar_eq_of_not_hasFVar}
Structural inductions, each mutual with its three list twins; the only arithmetic is the
cutoff bookkeeping at the variable, $\lambda$, \code{let}, case and fixpoint nodes.
\end{proof}
\end{lemma}

\section{Weakening, substitution, closing, universe instantiation}
\label{sec:erases:transport}

\begin{theorem}[Erasure commutes with de Bruijn weakening]
\label{thm:erases_shift}
\lean{LeanToLambdaBox.erases_shift}
\leanok
\uses{def:Erases, def:LBTerm-shift, def:lean4lean-grounding}
Along a weakening of the local context by de Bruijn entries, lifting the source matches
shifting the target: if $e$ erases to $t$, then the lifted source erases to the shifted
target. The Lean statement carries lean4lean's bound-variable lift witness, which fixes both
cutoffs, and asks only that the environment be ordered. This is the weakening half of the
[S, Sect. 7.3] and [L, Lemma 16] package, with the target side included, which the paper
statements leave implicit.
\begin{proof}
\leanok
\uses{def:Erases, lem:Erasable-weakN}
Induction on the derivation, generalising the target context and the cutoff. The \code{box}
arm moves its two witnesses by lean4lean's translation weakening and by \code{Erasable.weakN};
the variable arms move their lookups by the lift witness's own lookup lemma, whose index
convention both sides already use; the structural arms are congruences.
\end{proof}
\end{theorem}

\begin{theorem}[Erasure commutes with substitution]
\label{thm:erases_subst}
\lean{LeanToLambdaBox.erases_subst, LeanToLambdaBox.erases_subst_let,
LeanToLambdaBox.instN_toBVLift, LeanToLambdaBox.instLet_toBVLift}
\leanok
\uses{def:Erases, def:LBTerm-subst, def:lean4lean-grounding}
If the substitutee $e_0$ erases to $s'$, then substituting $e_0$ into $e$ at depth
$\mathit{dk}$ on the source matches substituting $s'$ into $t$ at depth $\mathit{dk}$ on the
target; and the same holds when the context entry instantiated away is a \code{let} entry
whose recorded value is the substitutee's translation. Eight lookup facts sit underneath, of
which the two cited say that an instantiation witness yields the de Bruijn weakening of the
substitutee's context into the instantiated one; the other six say that a bound variable below
the substituted index keeps its index, one above it drops by one, and a free variable's lookup
is untouched. The Lean statements additionally require the substitutee's translation and its
typing. These are the lemmas the $\beta$ and $\zeta$ arms of the simulation spend
(Chapter~\ref{chap:correctness}); [L, Lemma 16] is the informal statement.
\begin{proof}
\leanok
\uses{def:Erases, lem:Erasable-weakN, thm:erases_shift}
Induction on the derivation. The \code{box} arm is discharged by lean4lean's translation
substitution together with \code{Erasable.inst}. The variable arm at the substituted index is
the substitutee's own derivation lifted by Theorem~\ref{thm:erases_shift} along the derived
weakening; the other variable cases are the two lookup shifts; the binder arms push the
instantiation under the extended context. In the \code{let} variant the irrelevance witness
needs no move at all, since a \code{let} entry contributes nothing to the pure typing context.
\end{proof}
\end{theorem}

\begin{lemma}[Transport along a definitionally equal context]
\label{lem:Erases-defeqDFC_wt}
\lean{LeanToLambdaBox.Erases.defeqDFC_wt}
\leanok
\uses{def:Erases, def:lean4lean-grounding}
An erasure derivation transports along a definitionally equal local context, given a model
translation of the source in the original context and that context's well-formedness. It is
consumed by the $\zeta$ arm of the simulation, where a \code{let} entry's recorded value is
replaced by a convertible one.
\begin{proof}
\leanok
\uses{def:Erases, lem:Erasable-weakN, lem:Erasable-defeq, asm:lean4lean-trust}
Induction on the derivation. The \code{box} arm moves its witnesses by lean4lean's
context-transport lemma, by \code{Erasable.defeqDFC} and by Lemma~\ref{lem:Erasable-defeq}.
The binder arms must push a definitional equality for the entry they extend by, which is typed
at a sort and therefore needs the binder type's well-typedness --- data the \code{lam} rule
does not carry and the assumed translation does, which is exactly why the translation is a
premise. The proof routes through lean4lean's uniqueness of translations, and the measured
axiom footprint contains \code{sorryAx}.
\end{proof}
\end{lemma}

\begin{theorem}[Open, erase, close]
\label{thm:Erases-abstract}
\lean{LeanToLambdaBox.Erases.abstract, LeanToLambdaBox.Erases.uninstantiateN,
LeanToLambdaBox.Erases.uninstantiate, LeanToLambdaBox.closed_instantiate1'_fvar,
LeanToLambdaBox.abstract_find?_dk, LeanToLambdaBox.fvarId_beq_comm}
\leanok
\uses{def:Erases, def:toBvar, def:lean4lean-grounding}
Erasure commutes with closing a free variable back into a binder: if $e$ erases to $t$ in a
context whose entry at depth $\mathit{dk}$ is the free variable $v_0$, and $e$ has no loose
bound variable at or above $\mathit{dk}$, then the source closed at $v_0$ erases to
$\code{toBvar}\ v_0\ \mathit{dk}\ t$ in the context with that entry flipped to a de Bruijn
binder. The form the eraser bridge actually calls is the converse reading: if the binder body
\emph{opened} with a fresh $v_0$ erases to $t$, then the un-opened body erases to the closed
target --- in one word, \emph{open, erase, close} equals \emph{erase}. The Lean statements
carry lean4lean's context-abstraction witness, a closedness side condition, and, for the
derived form, freshness of $v_0$ in the body. This is the reconciliation for the shipping
eraser's open-recurse-abstract idiom, the one part of the metatheory that neither [S] nor [L]
needs, both working with de Bruijn terms throughout ([R, Sect. 4.7]).
\begin{proof}
\leanok
\uses{def:Erases, lem:toBvar_eq_of_not_hasFVar}
Induction on the derivation. The variable arms are the interesting ones: after the context
surgery the entry at depth $\mathit{dk}$ is a de Bruijn binder, and the free variable being
closed is exactly the one the index now denotes --- with a small commutativity step, because
the source-side abstraction tests $v_0$ against the bound variable while \code{toBvar}
tests the bound variable against $v_0$. The
closedness premise is needed because source abstraction shifts loose bound variables at or
above the insertion level while \code{toBvar} does not, and the two agree only when there is
no such variable. The two derived forms follow by rewriting with the freshness hypothesis.
\end{proof}
\end{theorem}

\begin{definition}[The \texttt{max}-free level fragment]
\label{def:NoMaxLevels}
\lean{LeanToLambdaBox.NoMaxLevel, LeanToLambdaBox.NoMaxLevels,
LeanToLambdaBox.noMaxLevels_toConstructor}
\leanok
\code{NoMaxLevel} holds of a universe level built without \code{max} and \code{imax}, and
$\code{NoMaxLevels}\ e$ says that every level $e$ mentions in its \code{sort} and \code{const}
nodes is such. The restriction exists because Lean's level substitution calls the
\emph{normalising} smart constructors for \code{max} and \code{imax}, so level substitution is
loose up to conversion at such a node; excluding those nodes makes the transport an equation.
A literal's constructor unfolding mentions only the levels $[]$ and $[\code{zero}]$, so it
lies in the fragment.
\end{definition}

\begin{lemma}[Strict level substitution]
\label{lem:substParams_strict}
\lean{LeanToLambdaBox.substParams_strict, LeanToLambdaBox.substParams_strict_list,
LeanToLambdaBox.TrExprS.instL_core, LeanToLambdaBox.TrExprS.instL_strict}
\leanok
\uses{def:NoMaxLevels, def:lean4lean-grounding}
On \code{max}-free levels, level substitution is strict: reading the substituted level yields
\emph{exactly} the instantiated model level, where the upstream statement concludes only
equivalence. Consequently, on the \code{max}-free fragment, level instantiation of a
translation is again a translation \emph{on the nose}, where lean4lean's own instantiation
lemma concludes only translation up to conversion. Because the transport is strict, the strict
version needs neither environment nor context well-formedness.
\begin{proof}
\leanok
\uses{def:NoMaxLevels, asm:lean-reflection-axioms}
Induction on the level, respectively on the translation; the excluded nodes are exactly those
at which the smart constructors could normalise. The measured axiom footprint contains
lean4lean's reflection axiom for level parameters.
\end{proof}
\end{lemma}

\begin{remark}
\code{TrExprS.instL\_core} and \code{TrExprS.instL\_strict} are \textbf{repository-local}
lemmas declared as \code{theorem TrExprS.\dots} inside \code{namespace LeanToLambdaBox}: their
real names carry that prefix and they are not lean4lean's. They are re-proofs of upstream
inductions under a stronger conclusion, hence proof scripts tied to the lean4lean pin. The
same warning applies to several names in Section~\ref{sec:erases:strengthen}.
\end{remark}

\begin{theorem}[Erasure transports along universe instantiation]
\label{thm:Erases-instL}
\lean{LeanToLambdaBox.Erases.instL, LeanToLambdaBox.Erases.instL_core,
LeanToLambdaBox.erases_instL_closed}
\leanok
\uses{def:Erases, def:NoMaxLevels}
On the \code{max}-free fragment, instantiating the source's universe parameters leaves its
\lbox{} image \textbf{literally unchanged}: if $e$ erases to $t$ at parameters $\mathit{ps}$,
then the source with those parameters instantiated erases to the same $t$ at the new scope, in
the instantiated context. That \lbox{} carries no universes is the whole content of the
theorem. The Lean statement asks that the substituted levels be readable in the target scope
and that the two lists have equal length.
\begin{proof}
\leanok
\uses{def:Erases, lem:substParams_strict, lem:Erasable-weakN, asm:lean-reflection-axioms}
Induction on the derivation with the target fixed. The \code{box}, \code{lam} and \code{letE}
arms move their translations by Lemma~\ref{lem:substParams_strict} and the irrelevance witness
by \code{Erasable.instL}; the \code{ctor}, \code{const} and \code{proj} arms are untouched,
their premises being about the environment rather than about levels. The measured axiom
footprint adds the \code{Lean.Expr} reflection axioms and two \code{bv\_decide} certificates.
\end{proof}
\end{theorem}

\begin{remark}
A small theorem with a large consequence: it is what makes the strong definition clause of
\ErasesEnv{} --- one \lbox{} body an erasure of the source body at \emph{every} scope and
instantiation --- satisfiable from a single derivation. The clause itself does not mention
\code{NoMaxLevels}, which this theorem requires; the gap is recorded in
Chapter~\ref{chap:trust}.
\end{remark}

\section{Free-variable transport and context uniformity}
\label{sec:erases:strengthen}

\begin{lemma}[Dropping an unused \texttt{let} entry]
\label{lem:Erases-thin_vlet}
\lean{LeanToLambdaBox.Erases.thin_vlet, LeanToLambdaBox.Erases.strengthen_vlet,
LeanToLambdaBox.ThinVLet, LeanToLambdaBox.ThinVLet.toCtx,
LeanToLambdaBox.ThinVLet.fvars_eq, LeanToLambdaBox.ThinVLet.find?,
LeanToLambdaBox.TrExprS.thin_vlet}
\leanok
\uses{def:Erases, def:lean4lean-grounding}
A derivation at a context containing an \textbf{unused} free-variable-tagged \code{let} entry
also holds with that entry dropped: same source, same target, and no well-formedness or
closedness premise needed. The context surgery \code{ThinVLet} describes exactly the contexts
reachable from such an entry by the binder rules, which cons only de Bruijn entries; it leaves
the pure typing context untouched, removes exactly the one free variable from the context's
variable list, and leaves lookup unchanged at every other variable. At depth zero, a
\code{let}'s value erased under a freshly opened binder can be read at the outer context,
since the value cannot mention that binder. The Lean statements carry a freshness side
condition on the source. The lemma is consumed by the eraser bridge
(Chapter~\ref{chap:refinement}).
\begin{proof}
\leanok
\uses{def:Erases}
Induction on the derivation; every lean4lean side premise transports on the nose, because the
lookup fact is an \emph{equality} and the model witness is unchanged --- which is what the
repository-local \code{TrExprS.thin\_vlet} establishes for the translation.
\end{proof}
\end{lemma}

\begin{lemma}[Free-variable weakening]
\label{lem:erases_weakFV}
\lean{LeanToLambdaBox.erases_weakFV, LeanToLambdaBox.TrExprS.weakFV_fvwf,
LeanToLambdaBox.TrExprS.weakFV'_fvwf, LeanToLambdaBox.VLCtx.FVWF.fvars_nodup,
LeanToLambdaBox.VLCtx.FVLift.find?_fvwf, LeanToLambdaBox.VLCtx.FVLift'.find?_fvwf}
\leanok
\uses{def:Erases, def:lean4lean-grounding}
A derivation at a context also holds at a context extended by free-variable entries. The Lean
statement asks only for a \emph{free-variable} well-formedness of the larger context ---
distinctness of its variable list --- where the upstream weakening asks for full
well-formedness; that relaxation is why five of the six cited names are repository-local
re-proofs of upstream inductions, declared under upstream-looking namespaces and tied to the
lean4lean pin.
\begin{proof}
\leanok
\uses{def:Erases, lem:Erasable-weakN}
Induction on the derivation, with the lookup arms discharged by the re-proved transport, the
\code{box} arm by the re-proved translation weakening together with \code{Erasable.weakN}, and
the rest by congruence. The re-proofs follow upstream's own scripts, which touch the
well-formedness hypothesis only through the tail and through distinctness of the variable
list, both of which the weaker predicate supplies.
\end{proof}
\end{lemma}

\begin{lemma}[Weakening with no hypothesis on the context]
\label{lem:erases_weakFV_nofvars}
\lean{LeanToLambdaBox.erases_weakFV_nofvars, LeanToLambdaBox.TrExprS.weakFV_nofvars,
LeanToLambdaBox.TrExprS.weakFV'_nofvars, LeanToLambdaBox.VLCtx.FVLift.find?_inl,
LeanToLambdaBox.VLCtx.FVLift'.find?_inl}
\leanok
\uses{def:Erases, def:lean4lean-grounding}
The same weakening with the target context's well-formedness premise removed
\textbf{entirely}, paid for by free-variable-freeness of the source: a bound-variable lookup
never reaches the branch that needs distinctness --- that branch exists only to refute a
shadowing free variable --- and the free-variable rule cannot fire at all. This is the
ingredient that makes weakening to an \emph{arbitrary} context possible.
\begin{proof}
\leanok
\uses{def:Erases, lem:erases_weakFV, lem:Erasable-weakN}
As for Lemma~\ref{lem:erases_weakFV}, with the two lookup lemmas replaced by their
bound-variable-only variants and the free-variable arm refuted by the freeness hypothesis.
\end{proof}
\end{lemma}

\begin{theorem}[Unrestricted weakening]
\label{thm:erases_weak_any}
\lean{LeanToLambdaBox.erases_weak_any}
\leanok
\uses{def:Erases, def:LBClosed, def:lean4lean-grounding}
For a closed, free-variable-free source erasing to a closed \lbox{} term, a derivation at the
empty context holds at \emph{every} context, with no hypothesis on that context whatsoever.
This is the form in which a closed top-level body's erasure is replayed wherever it is needed:
the weakening leg of context uniformity.
\begin{proof}
\leanok
\uses{thm:erases_shift, lem:erases_weakFV_nofvars}
Induction on the context, adding one entry at a time. A de Bruijn entry goes through
Theorem~\ref{thm:erases_shift}, whose two lifts are the identity by closedness on both sides;
a variable entry --- possibly shadowing --- goes through
Lemma~\ref{lem:erases_weakFV_nofvars}, which asks nothing of the context.
\end{proof}
\end{theorem}

\begin{definition}[A commissioned premise: strengthening irrelevance]
\label{def:ErasableStrengthen}
\lean{LeanToLambdaBox.ErasableStrengthen, LeanToLambdaBox.erasableStrengthen_liftN_zero,
LeanToLambdaBox.Ctx.LiftN.zero_eq}
\leanok
\uses{def:Erasable}
$\code{ErasableStrengthen}\ \mathit{env}\ \mathit{Us}$ is an \textbf{inverse of}
\code{Erasable.weakN}: erasability of a lifted term in a larger context implies erasability of
the term in the smaller one. \code{Erasable} ships with weakening, instantiation, conversion
and context transport and with no inverse of weakening, and the pinned lean4lean has none for
its typing judgement either. The obligation is therefore \textbf{named as a premise}, visible
in the statements that need it, rather than assumed as an axiom. Its hard half is the arity
disjunct, which needs a $\Pi$ inversion through a lift. A lift by zero satisfies the premise as
the identity, which pins the quantifier structure and shows it is not vacuously false.
\end{definition}

\begin{remark}
This is \emph{not} a field of any capstone hypothesis bundle, so it is stated here as a
definition rather than as an assumption of Chapter~\ref{chap:trust}, which classifies it as a
named premise and never as an axiom. The non-vacuity witness is weak: it holds only at a zero
lift, where the statement degenerates, and gives no evidence at a genuine lift.
\code{Ctx.LiftN.zero\_eq} is likewise repository-local under an upstream-looking namespace.
\end{remark}

\begin{definition}[Two scope conditions for strengthening]
\label{def:NoProj}
\lean{LeanToLambdaBox.NoProj, LeanToLambdaBox.NoProjBinders,
LeanToLambdaBox.NoProj.toIsUnique, LeanToLambdaBox.NoProj.toNoProjBinders,
LeanToLambdaBox.NoProj.toConstructor, LeanToLambdaBox.NoProjBinders.toConstructor,
LeanToLambdaBox.NoProj.natLitToConstructor, LeanToLambdaBox.NoProj.strLitToConstructor}
\leanok
\uses{def:lean4lean-grounding}
$\code{NoProj}\ e$ is projection-freeness at every subterm, \textbf{including a \code{let}
binder's type annotation}; it is strictly stronger than lean4lean's own uniqueness scope
condition, which omits that annotation because the translation of a \code{let} does not depend
on it --- but the \code{letE} rule of \Erases{} records that type in the body's context entry,
so the strengthening must pin it. $\code{NoProjBinders}\ e$ is the weaker condition that asks
projection-freeness only where the strengthening actually spends uniqueness: at $\lambda$ and
$\forall$ binder types and at a \code{let}'s type and bound value. Projections are permitted
everywhere computational --- $\lambda$ and \code{let} bodies, application heads and arguments,
and under a projection --- these being exactly the positions at which \Erases{} records a model
witness that a context lift must match on the nose. The binder clauses cannot be dropped:
equational uniqueness at a projection is \textbf{false}, not merely unproved, because
lean4lean's projection premise pins parameters only up to definitional equality. Every
literal's one-step unfolding satisfies both conditions.
\end{definition}

\begin{lemma}[The two conditions are different]
\label{lem:noProjBinders_ofNatBody}
\lean{LeanToLambdaBox.ofNatBody, LeanToLambdaBox.noProjBinders_ofNatBody,
LeanToLambdaBox.noProj_ofNatBody_refuted}
\leanok
\uses{def:NoProj}
The prepared body of a class method, $\mathrm{fun}\ (\alpha : \mathrm{Type})\ (x :
\mathrm{Nat})\ (\mathit{self} : \code{OfNat}\ \alpha\ x) \Rightarrow \mathit{self}.1$,
satisfies the weaker condition --- its three binder types are projection-free and its
projection sits in the body --- and provably \textbf{fails} the strong one. That is what the
relaxation buys, and why the second strengthening induction is the usable one: the corpus
contains class-method projections.
\begin{proof}
\leanok
\uses{def:NoProj}
Both by unfolding the two predicates at the literal term.
\end{proof}
\end{lemma}

\begin{theorem}[Strengthening at depth]
\label{thm:Erases-strengthen_fvlift}
\lean{LeanToLambdaBox.Erases.strengthen_fvlift,
LeanToLambdaBox.Erases.strengthen_fvlift_binders}
\leanok
\uses{def:Erases, def:ErasableStrengthen, def:NoProj, def:lean4lean-grounding}
A derivation at a free-variable extension of a context also holds at the smaller context,
provided the source satisfies the scope condition and is \emph{already translatable} at the
small context, and provided the commissioned premise of
Definition~\ref{def:ErasableStrengthen}. There are two versions, one at the strong scope
condition and one at the weaker one; the second pays for the relaxation in the \code{box} arm,
where lean4lean's \emph{conversion} uniqueness plus Lemma~\ref{lem:Erasable-defeq} replace the
equational uniqueness, an upgrade that moves two premises --- from free-variable-only to full
context well-formedness, and from an ordered to a well-formed environment. Both are proved
\emph{modulo the named premise}, and neither is in the capstone's import closure.
\begin{proof}
\leanok
\uses{def:Erases, lem:erases_weakFV, lem:Erasable-defeq, def:ErasableStrengthen,
asm:lean4lean-trust}
Induction on the derivation. The assumed small-context translation is the engine: pushed out to
the large context by the re-proved translation weakening and identified with the derivation's
own witness by uniqueness, it turns every translation-bearing arm into a rewrite, and it
supplies the lookup premises of the variable rules at the small context, which weakening got
for free and strengthening cannot. The commissioned premise is spent exactly once, in the
\code{box} arm. The route through lean4lean's inverse-weakening lemma is unavailable: it
recovers a small-context translation only up to definitional equality, which is survivable in
\code{box} and fatal in \code{lam}, where the body's induction hypothesis demands an equation
between binder types. Of the two versions, the measured axiom footprint of the weak-scope one
contains \code{sorryAx}, through the conversion uniqueness; that of the strong-scope one does
not.
\end{proof}
\end{theorem}

\begin{theorem}[Context uniformity]
\label{thm:erases_uniform_closed}
\lean{LeanToLambdaBox.erases_strengthen_closed, LeanToLambdaBox.erases_uniform_closed,
LeanToLambdaBox.erases_uniform_of_nil}
\leanok
\uses{def:Erases, def:ErasableStrengthen, def:NoProj, def:LBClosed,
def:lean4lean-grounding}
For a closed, free-variable-free source whose binders are projection-free and whose image is
closed, the erasure does not depend on the local context it was produced at: a derivation at
one context holds at every other. The route is to strengthen to the empty context and then
weaken to an arbitrary one. Strengthening out of the empty context uses the only shape
available --- a lift out of the empty context can only skip variable entries --- which is
exactly the eraser's situation, a top-level body erased under a context of opened free
variables; closedness and freeness of the source are consequences of the assumed translation,
not premises. The one-sided corollary at the empty context needs \textbf{no} commissioned
premise: only the two-sided transport, which must also come back down from the call site's
context, buys the obligation.
\begin{proof}
\leanok
\uses{thm:Erases-strengthen_fvlift, thm:erases_weak_any, asm:lean4lean-trust}
Compose Theorem~\ref{thm:Erases-strengthen_fvlift}, in its weak-scope form, with
Theorem~\ref{thm:erases_weak_any}. The second leg cannot be the free-variable weakening of
Lemma~\ref{lem:erases_weakFV}, because the target context is arbitrary --- de Bruijn entries
and shadowing variable entries included --- which is what breaks its well-formedness premise.
The measured axiom footprint contains \code{sorryAx}, inherited through the first leg.
\end{proof}
\end{theorem}

\begin{remark}
An honesty flag. The intended consumer of this theorem is a bridge step at which the erasure
of a constant body produced under one context must be read under another. But the module
\code{ErasesUniform.lean} is outside the capstone's import closure and currently has no
consumer at all; \code{doc/coverage.md} counts its nineteen declarations in the
dead-declaration budget. The theorem is therefore not a step in the main line of argument.
\end{remark}

\section{Environment erasure}
\label{sec:erases:env}

\begin{definition}[Environment erasure]
\label{def:ErasesEnv}
\lean{LeanToLambdaBox.ErasesEnv, LeanToLambdaBox.ErasesEnv.mk}
\leanok
\uses{def:Erases, def:ConstOrigin, def:IndInfo, def:IndDeclOf, def:InformativeInd,
def:ReachableFrom, def:ElimDecl, def:CasesOnShape, def:isCasesOnName, def:Kername,
def:LBTerm-envLookup, def:GlobalDeclarations}
$\ErasesEnv\ \mathit{env}\ \mathit{bo}\ \Gamma_{\mathrm{spec}}\ t$ says that the \lbox{}
environment $\Gamma_{\mathrm{spec}}$ is a correct, \textbf{dependency-selective} erasure of
the source environment, \emph{as seen from the program} $t$. Here $\mathit{bo}$ is the
\textbf{compiler body table}, the table the source $\delta$ rule reads, so that source and
target unfold the same term. This is [S]'s \code{erases\_deps} (p. 8:64) rather than the
pointwise erasure of [S, Table 2]: every clause is triggered by the program's reachability of
a kername, and every clause is stated in \code{erases\_deps}' own direction --- the source
declares $X$, the target declares $X'$, and $X'$ is an erasure of $X$. There are seven
clauses:

\medskip
\noindent\begin{tabular}{lp{9.9cm}}
\code{keys} & the environment's keys are distinct (the \code{fresh\_global} half of [S]'s
well-formedness, so that a shadowing entry cannot retarget a pass) \\
\code{deps} & every kername the program reaches is declared \\
\code{tabled} & a constant the compiler table defines is a plain constant of the model; there
is no reachability trigger, because the reading it refutes is at a constant for which the
erasure emits no key \\
\code{defns} & every reached, tabled constant's entry is a declaration with a body, and that
body is an erasure of the compiler body \emph{at every level scope and instantiation} \\
\code{axioms} & a reached, body-less, non-eliminator constant is declared body-less, without
which the environment could hold a junk body where the source cannot step and the target
$\delta$-unfolds \\
\code{blocks} & a reached inductive block is declared in the environment's own declaration
list, and its \lbox{} entry carries the arity and propositionality data the target reads \\
\code{elims} & a reached \code{casesOn} of a \emph{relevant} inductive has the eliminator's
canonical body, at the segmentation the block fixes \\
\end{tabular}
\medskip

There is deliberately \emph{no} constructor clause: a constructor constant has no entry of its
own, and what a constructor value's arity is read off is the block, which the \code{blocks}
clause supplies at the block kername a constructor node already reaches.
\end{definition}

\begin{remark}
The \code{defns} clause is very strong: it quantifies over \emph{all} level scopes and
instantiations, including ill-matched ones. It is satisfiable only because \lbox{} carries no
universes (Theorem~\ref{thm:Erases-instL}), and even that theorem needs the \code{max}-free
fragment the clause does not mention. This clause is one of the two fields of the capstone's
bridge bundle that no theorem yet produces for a run of the shipping eraser
(Chapter~\ref{chap:trust}). Formally \ErasesEnv{} is a one-constructor \code{inductive}, not a
\code{structure}, so its clauses are reached through hand-written accessors rather than
generated projections.
\end{remark}

\begin{lemma}[The seven accessors]
\label{lem:ErasesEnv-keys}
\lean{LeanToLambdaBox.ErasesEnv.keys, LeanToLambdaBox.ErasesEnv.deps,
LeanToLambdaBox.ErasesEnv.tabled, LeanToLambdaBox.ErasesEnv.defns,
LeanToLambdaBox.ErasesEnv.axioms, LeanToLambdaBox.ErasesEnv.blocks,
LeanToLambdaBox.ErasesEnv.elims}
\leanok
\uses{def:ErasesEnv}
The seven clauses are available as separate lemmas, and these are the names the arms of the
simulation cite (Chapter~\ref{chap:correctness}). Two additionally take their trigger premises
as explicit arguments: the block accessor takes the block data and the reachability of the
block kername, and the eliminator accessor takes the eliminator's shape, the relevance of its
inductive, its constant origin, and the reachability of its kername.
\begin{proof}
\leanok
\uses{def:ErasesEnv}
Each is a case analysis on the single constructor.
\end{proof}
\end{lemma}

\begin{definition}[Coverage of one inductive]
\label{def:IndCovered}
\lean{LeanToLambdaBox.IndCovered}
\leanok
\uses{def:IndInfo, def:IndDeclOf, def:InformativeInd, def:ConstOrigin, def:ElimDecl,
def:CasesOnShape}
$\Gamma_{\mathrm{spec}}$ \textbf{covers} the inductive $n$ when it declares $n$'s block with
the numbers the target semantics reads and --- when $n$ is relevant --- declares its
\code{casesOn} eliminator at the segmentation the block fixes. These are \ErasesEnv{}'s two
inductive clauses at a single name, which is the granularity at which one registration step of
the shipping eraser establishes them. Formally it is a two-field structure.
\end{definition}

\begin{definition}[What a specification environment says about the source]
\label{def:SpecContent}
\lean{LeanToLambdaBox.SpecContent}
\leanok
\uses{def:Erases, def:IndCovered, def:ConstOrigin, def:IndInfo, def:CasesOnShape,
def:isCasesOnName}
$\code{SpecContent}\ \mathit{env}\ \mathit{bo}\ \Gamma_{\mathrm{spec}}$ bundles \ErasesEnv{}'s
five source-facing clauses with \textbf{the entry's presence} in place of \textbf{the
program's reachability}. Being program-independent, the bundle is fixed for a run, which is
what makes it maintainable as a loop invariant by the eraser's registration path. The
\code{tabled} clause is not among the five, because it mentions no environment at all. One
difference matters: the definition field records \textbf{one} level scope --- what a run
actually erases at --- where \ErasesEnv{}'s asks for every one; and the two inductive fields
conclude coverage in the sense of Definition~\ref{def:IndCovered}. The bundle is consumed by
the registry invariant of Chapter~\ref{chap:coldstart}.
\end{definition}

\begin{lemma}[Reading the content at a program]
\label{lem:SpecContent-erasesEnv}
\lean{LeanToLambdaBox.SpecContent.erasesEnv}
\leanok
\uses{def:SpecContent, def:ErasesEnv}
A specification environment's content, read at a program, is an environment erasure for that
program: the clauses whose trigger is the program's reachability get it through a
dependency-closure premise. Two further premises are supplied separately --- the all-scopes form of
the definition clause, because the content bundle records one scope, and the \code{tabled}
clause, because it is not a fact about the environment at all. Despite its name the lemma does
\emph{not} use the bundle's own definition field; it is better described as the transport of
the four remaining fields plus the three premises just listed (dependency closure, the
all-scopes \code{defns} form, and \code{tabled}).
\begin{proof}
\leanok
\uses{def:SpecContent, def:ErasesEnv, def:IndCovered}
Apply the single constructor, feeding each reachability trigger through the dependency premise
and unpacking the two inductive fields' coverage conclusions into the block and eliminator
clauses.
\end{proof}
\end{lemma}

\begin{definition}[Well-formedness of a specification environment]
\label{def:LBWfSpec}
\lean{LeanToLambdaBox.LBWfSpec}
\leanok
\uses{def:ElimDecl, def:GlobalDeclarations}
$\code{LBWfSpec}\ \Gamma_{\mathrm{spec}}$ is distinct keys together with closed bodies --- the
closedness on which the pass relation's shift and substitution commutations run. It appears in
the capstone's environment existential (Chapter~\ref{chap:capstone}).
\end{definition}

\begin{definition}[The emitted environment as a lowered, pruned image]
\label{def:LowerEnv}
\lean{LeanToLambdaBox.LowerEnv}
\leanok
\uses{def:Lower, def:LowerFix, def:ElimDecl, def:LBTerm-envLookup, def:GlobalDeclarations}
$\LowerEnv\ \Gamma_{\mathrm{spec}}\ \Gamma$ says that the emitted \lbox{} environment $\Gamma$
is the lowered, pruned image of the specification environment $\Gamma_{\mathrm{spec}}$, which
is itself well formed for the pass. Its eight fields: the emitted keys are distinct; a body
that both environments declare at the same key is related by the pass relation, or is one
member of a lowered mutual-fixpoint block; every definition that is not a runtime key survives
the pruning as a definition --- without which a pruned definition would satisfy the previous
field vacuously while the target is stuck at its constant node; a body-less constant stays
body-less or is pruned; inductive blocks are carried over unchanged, since the target's
constructor-arity and propositionality queries read the specification's numbers; pruning only
removes; and both environments' bodies are closed.
\end{definition}

\begin{remark}
This is the interface between the erasure specification and the pass layer of
Chapter~\ref{chap:lowering}; together with \ErasesEnv{} it forms the two fields of the
capstone's bridge bundle. The declaration lives in \code{ErasesEnv.lean}, which is why it is
defined in this chapter and used by Chapters~\ref{chap:lowering} and~\ref{chap:capstone}.
\end{remark}

\begin{lemma}[A lowered environment exists]
\label{lem:lowerEnv_idEnv}
\lean{LeanToLambdaBox.idEnv, LeanToLambdaBox.lowerEnv_idEnv}
\leanok
\uses{def:LowerEnv, def:Kername}
The smallest environment carrying a real pass step --- one definition whose body is a closed
identity $\lambda$, kept by the pruning --- is its own lowered image, with all eight fields
discharged.
\begin{proof}
\leanok
\uses{def:LowerEnv, def:Lower}
Decide the key list; the single body is related to itself by the pass's $\lambda$ congruence;
the remaining fields follow by unfolding the one-entry lookup.
\end{proof}
\end{lemma}

\begin{remark}
An honest caveat: on a one-definition environment the body-less, inductive and totality fields
hold vacuously --- there is no body-less constant, no inductive entry and no runtime key.
\end{remark}

\begin{definition}[A four-entry specification environment]
\label{def:demoEnv}
\lean{LeanToLambdaBox.demoInd, LeanToLambdaBox.demoCases, LeanToLambdaBox.demoDef,
LeanToLambdaBox.demoAx, LeanToLambdaBox.demoEnv, LeanToLambdaBox.demoProg,
LeanToLambdaBox.DemoSource}
\leanok
\uses{def:IndInfo, def:IndDeclOf, def:ConstOrigin, def:Erases, def:mkElimBody,
def:CasesOnShape, def:OneInductiveBody, def:Kername}
A literal environment exercising every kind of entry a specification environment holds: an
inductive block, its \code{casesOn} eliminator, a recursive definition and a body-less axiom,
with the block and the eliminator together, which is the configuration the two inductive
clauses are read at. Its program is the recursive definition itself, whose body calls the
eliminator on a recursive call --- the recursion being a constant node here, since a fixpoint
node is what the \emph{lowering} makes of it. \code{DemoSource} bundles the eleven source-side
facts the entries are read against: which constants the compiler table defines, the fixture's
own block and definition, block and eliminator uniqueness, and four kername-separation
clauses, the last needed because the name-to-kername map is \textbf{not} injective.
\end{definition}

\begin{remark}
The block kername is spelled as the literal the minting function evaluates to, because
\code{Name.toString} does not reduce definitionally --- a small but real obstacle, since
reachability through a case node computes with that kername.
\end{remark}

\begin{theorem}[The seven clauses are jointly satisfiable]
\label{thm:demoEnv_erasesEnv}
\lean{LeanToLambdaBox.demoEnv_erasesEnv, LeanToLambdaBox.demoEnv_keys,
LeanToLambdaBox.demoEnv_indCovered, LeanToLambdaBox.demoEnv_specContent,
LeanToLambdaBox.demoEnv_deps, LeanToLambdaBox.demoEnv_defns}
\leanok
\uses{def:ErasesEnv, def:demoEnv}
Relative to the source-side bundle, the four-entry fixture \textbf{is} a specification
environment for its own program: its four keys are distinct and any key it answers is one of
them; it covers its own inductive, with the block entry and the canonical eliminator body at
the numbers the source side reads; its entries satisfy the program-independent content bundle;
its program reaches exactly the four declared kernames --- the axiom, the block the
eliminator's case node reads, the eliminator, and the definition itself; and every reached,
tabled constant's body erases to the entry the fixture holds, at every level scope.
\begin{proof}
\leanok
\uses{lem:SpecContent-erasesEnv, def:SpecContent, def:IndCovered, def:demoEnv,
def:ReachableFrom}
Compute the reachable set of the program with the reachability function and the four lookup
equations. Establish the content bundle field by field from the source-side hypotheses, using
the fixture's key-case analysis to identify which entry each query hits. Then apply
Lemma~\ref{lem:SpecContent-erasesEnv}, supplying the all-scopes definition clause from the
bundle's stability field.
\end{proof}
\end{theorem}

\begin{remark}
This is \emph{conditional} non-vacuity, and must be read as such: the eleven-field source
bundle is a set of \emph{hypotheses} about a source environment one cannot hand-build, so the
fixture shows that the seven clauses are jointly satisfiable \emph{given} such a source, not
outright --- in contrast to Theorem~\ref{thm:erases_ctor_fires}, which is unconditional.
\end{remark}

\section*{Deviations from the paper, and caveats}

The relation of this chapter departs from [S, Fig. 18] in six places, each departure forced
rather than chosen. Constructors are in \emph{applied form}: the \code{ctor} rule fires at the
bare head and the constructor node carries a literally empty argument list, arguments arriving
through \code{app} (cf.\ shipping finding F-ETA2). There is no case rule and no fixpoint rule, because
\code{Lean.Expr} has neither node; both are declaration-level phenomena and belong to \Lower{}
(Chapter~\ref{chap:lowering}). Two arms are added that [S] has no need of, \code{mdata} and
\code{lit}. The \code{proj} rule carries a relevance premise that [S] gets from Rocq's typing,
and Theorem~\ref{thm:erases_proj_needs_informative} shows that dropping it would make the
corresponding arm of the simulation refutable rather than merely unprovable (design tag N18).
The single de Bruijn variable rule becomes two, because the shipping eraser works locally
nameless with freshly opened free variables ([R, Sect. 4.7]); the reconciliation is
Theorem~\ref{thm:Erases-abstract}, which neither [S] nor [L] needs. And \ErasesEnv{} is [S]'s
dependency-selective \code{erases\_deps}, not the pointwise environment erasure, with a
\code{defns} clause quantifying over \emph{every} level scope.

Three caveats travel with the chapter. First, the two environment relations defined here,
\ErasesEnv{} and \LowerEnv{}, are exactly the two fields of the capstone's bridge bundle, and
no theorem in the development yet produces that bundle for a run of the shipping eraser
(Chapter~\ref{chap:trust}). Second, the strong \code{defns} clause is satisfiable only through
Theorem~\ref{thm:Erases-instL}, which requires the \code{max}-free fragment that the clause
itself does not mention. Third, \code{ErasesUniform.lean} --- the commissioned premise of
Definition~\ref{def:ErasableStrengthen}, the two scope conditions, and the uniformity theorem
--- lies outside the capstone's import closure and has no consumer; it documents a route that
the closing round did not schedule.
